\scanpage{185}
avec \deleted{\unsure{une}}
\[
 (1\wedge x\wedge y\wedge z)\otimes e^{-1}=2\varepsilon_0\qquad
 \text{où}\quad\varepsilon_{\unsure{0|-}}=\unsure{[|(}\alpha\wedge\alpha'\wedge x_0\wedge y_0)\underset{\in\pm1}{\unsure{\otimes}e^{-1}}
\]

le signe dépend de la convention d'orientation (qu'il faudra préciser le moment venu !).

D'autre part
\[
 \deleted{\operatorname{Tr}xyz=1}\qquad \tau-\tau'=\operatorname{Tr}\underbrace{[x,y]}_z.z=\operatorname{Tr}z^2=\operatorname{Tr}1=2
\]

d'où
\[
\begin{aligned}
\deleted{\unsure{1|2}}2&=2\varepsilon\varepsilon_0\qquad\text{i.e.}\\
\varepsilon&=\deleted{\tfrac12}\varepsilon_0\ \deleted{\unsure{\varepsilon}}
\end{aligned}
\]

d'où
\[
 \boxed{\tau-\tau'=\deleted{\tfrac12}\varepsilon_0\underbrace{(1\wedge x\wedge y\wedge z)\otimes e^{-1}}_{\Delta\ \mathrm{déf}}}
 \qquad\text{i.e.}\quad 2(\tau-\tau')=\varepsilon_0\Delta
\]
(a un sens si $2$ inv. dans l'anneau de base $k$ où on travaille).

On en tire
\[
 4\tau\tau'=(\tau+\tau')^2-(\tau-\tau')^2=\Sigma^2-\deleted{\tfrac14}\Delta^2
\]

À vérifier (dans $\mathbb Q[[\deleted{b,b',b''}\unsure{X_{ij}|x_{ij}},y_{ij},\unsure{Z_{ij}|z_{ij}}]]$ \quad $1\leq i,j\leq2$)
\[
 \Sigma^2-\deleted{\tfrac14}\Delta^2\in4\mathbb Z[b,b',b'',t,t',t'']
\]
\margin{?}
\deleted{i.e.}
\[
 \deleted{4\Sigma^2-\Delta^2\in16\mathbb Z[b,b',b'',t,t',t'']}
\]

d'où l'expression
\[
 \tau\tau'=\Pi(b,b',b'',t,t',t'')\quad\in\mathbb Z[b,b',b'',\unsure{t|\psi},t',t'']
\]
expression polynomiale de poids total $6$ (les $b$ étant de
% archive sheet number: 184
